
\documentclass[12pt,a4paper]{article}
\usepackage{ctex}
\usepackage{geometry}
\usepackage{graphicx}
\usepackage{setspace}
\usepackage{indentfirst}
\usepackage{titlesec}
\usepackage{fancyhdr}
\usepackage{footmisc}
\usepackage{hyperref}

\geometry{left=2.5cm,right=2.5cm,top=2.5cm,bottom=2.5cm}
\setlength{\parindent}{2em}
\setlength{\headheight}{15pt}
\linespread{1.5}

\pagestyle{fancy}
\fancyhf{}
\fancyhead[C]{哈尔滨工业大学（深圳）马克思主义学院课程论文}
\fancyfoot[C]{\thepage}

\titleformat{\section}{\songti\zihao{4}}{\thesection}{1em}{}
\titleformat{\subsection}{\heiti\zihao{-4}}{\thesubsection}{1em}{}

\begin{document}

\begin{titlepage}
\thispagestyle{empty}
\begin{center}
\includegraphics[width=8cm]{校徽}

\vspace{0.3cm}
{\heiti\zihao{2} 马克思主义学院}

\vspace{0.5cm}
{\heiti\zihao{3} 2025年春季学期}

\vspace{0.6cm}
{\heiti\zihao{1} 课程论文}

\vspace{1.2cm}

\begin{tabular}{p{4cm}p{9cm}}
课程名称： & 毛泽东思想和中国特色社会主义理论体系概论\\
论文题目： & 邓小平理论对中国特色社会主义现代化建设的开创性贡献及时代价值\\
学生所在院（系）： & 机电工程与自动化学院\\
专业： & 电气工程及其自动化\\
学号： & 2024312506\\
姓名： & 余泓彬\\
上交日期： & 2026-06-11\\
考核结果： & \\
审阅人： & \\
\end{tabular}

\vfill
2026年
\end{center}
\end{titlepage}

\newpage
\begin{center}
{\heiti\zihao{3} 摘要}
\end{center}

邓小平理论是马克思主义中国化的重要成果，是中国特色社会主义理论体系的重要组成部分。本文围绕邓小平理论形成的历史背景、核心内容、对中国特色社会主义现代化建设的开创性贡献以及新时代的现实价值展开分析。研究认为，邓小平理论科学回答了“什么是社会主义、怎样建设社会主义”的重大问题，为改革开放和社会主义现代化建设提供了理论指导，并深刻影响了中国式现代化的发展道路。

\noindent\textbf{关键词：}邓小平理论；改革开放；中国特色社会主义；现代化建设；中国式现代化

\newpage

\begin{center}
{\heiti\zihao{-3} 邓小平理论对中国特色社会主义现代化建设的开创性贡献及时代价值}
\end{center}

\section{邓小平理论产生的时代背景}

20世纪70年代末，中国进入改革开放和社会主义现代化建设新时期。党的十一届三中全会实现了党和国家工作重心的历史性转移，为邓小平理论的形成奠定了现实基础。\footnote{《邓小平文选》第三卷，北京：人民出版社，1993年。}

新中国成立后，中国共产党带领人民建立了社会主义基本制度，但在社会主义建设过程中也经历了曲折探索。改革开放前，经济发展水平较低，人民生活改善速度有限。面对新的国内外形势，中国迫切需要寻找符合国情的发展道路。邓小平坚持实事求是，提出把马克思主义基本原理同中国具体实际相结合，推动理论创新与实践创新。

同时，世界科技革命和经济全球化趋势不断发展。国际竞争日益激烈，现代化建设任务更加紧迫。邓小平敏锐把握时代潮流，提出发展才是硬道理，强调必须集中力量发展社会生产力。

\section{邓小平理论的核心内容}

\subsection{解放思想、实事求是}

解放思想、实事求是是邓小平理论的重要思想基础。邓小平强调实践是检验真理的唯一标准，反对教条主义和经验主义，坚持从实际出发研究问题。改革开放的成功实践表明，只有不断解放思想，才能适应时代变化和社会发展要求。

\subsection{社会主义本质理论}

邓小平指出，社会主义的本质是解放生产力、发展生产力，消灭剥削，消除两极分化，最终达到共同富裕。\footnote{《邓小平文选》第三卷，第373页。}

这一论断突破了传统认识的局限，把生产力发展和共同富裕统一起来，为社会主义现代化建设指明了方向。社会主义不仅要坚持公有制和人民当家作主，更要不断提高社会生产力水平，使全体人民共享发展成果。

\subsection{改革开放理论}

改革开放是决定当代中国命运的关键一招。邓小平认为，改革是社会主义制度的自我完善和发展，对外开放是实现现代化建设的重要条件。通过改革开放，中国逐步建立社会主义市场经济体制，大幅提升了经济活力和国际竞争力。

\section{邓小平理论对中国特色社会主义现代化建设的开创性贡献}

\subsection{开辟中国特色社会主义道路}

邓小平坚持把马克思主义基本原理同中国具体实际结合起来，明确提出建设有中国特色社会主义。这一重大理论创新回答了中国如何建设社会主义、如何实现现代化的问题，成功开辟了中国特色社会主义道路。

中国特色社会主义道路的形成，不仅推动中国经济快速发展，也增强了国家治理能力和社会稳定水平。实践证明，这条道路符合中国国情，具有强大生命力。

\subsection{确立以经济建设为中心的发展战略}

邓小平提出“发展才是硬道理”，强调把经济建设作为党和国家工作的中心任务。通过实施改革开放政策，中国经济持续增长，人民生活水平显著提高。

以经济建设为中心的发展战略不仅推动工业化和城镇化进程，也为教育、科技、文化等事业发展提供了物质基础。事实证明，发展是解决中国一切问题的基础和关键。

\subsection{推动社会主义市场经济体制建立}

长期以来，人们普遍认为市场经济属于资本主义。邓小平突破传统观念束缚，指出计划和市场都是经济手段，不是区分社会主义与资本主义的本质标准。

在这一思想指导下，中国逐步建立社会主义市场经济体制，实现市场机制与社会主义制度优势相结合。社会主义市场经济既发挥市场配置资源的效率优势，又坚持共同富裕的发展方向，为经济持续增长提供了制度保障。\footnote{党的十四大提出建立社会主义市场经济体制目标。}

改革开放以来，中国经济总量实现历史性跨越，基础设施建设、科技创新能力和人民生活水平显著提高。这些成就与社会主义市场经济体制的建立密切相关。

\section{邓小平理论的时代价值}

进入新时代，中国式现代化全面推进。虽然时代条件发生了深刻变化，但邓小平理论仍然具有重要指导意义。

首先，坚持改革开放仍然是推动国家发展的关键动力。面对复杂多变的国际环境，中国必须继续深化改革、扩大开放，不断增强创新能力和国际竞争力。

其次，坚持发展生产力仍然是现代化建设的重要任务。高质量发展需要科技创新、产业升级和人才培养，而这些都与邓小平关于发展的思想一脉相承。

再次，坚持共同富裕体现社会主义本质要求。中国式现代化是全体人民共同富裕的现代化。新时代推进共同富裕，既是对邓小平社会主义本质理论的继承，也是中国特色社会主义发展的必然要求。

最后，坚持实事求是的方法论具有长期价值。面对新情况和新问题，必须坚持从实际出发，不断推进理论创新和制度创新。

\section{结论与展望}

邓小平理论是中国特色社会主义理论体系的重要组成部分，是改革开放和社会主义现代化建设的重要理论基础。它科学回答了什么是社会主义、怎样建设社会主义这一重大课题，成功开创了中国特色社会主义道路。

新时代新征程上，坚持和发展邓小平理论，对于推进中国式现代化、实现中华民族伟大复兴仍然具有重要意义。未来应继续坚持改革开放，坚持以人民为中心的发展思想，不断推进社会主义现代化建设。

\newpage
\begin{center}
{\heiti\zihao{3} 参考文献}
\end{center}

\begin{enumerate}
\item 邓小平.《邓小平文选》第三卷[M]. 北京：人民出版社，1993.
\item 邓小平.《邓小平文选》第二卷[M]. 北京：人民出版社，1994.
\item 王怀超. 当代中国社会主义改革的基本经验与历史启示[J]. 科学社会主义，2018(5):12--18.
\item 习近平.《习近平谈治国理政》第四卷[M]. 北京：外文出版社，2022.
\item 金冲及.《二十世纪中国史纲》[M]. 北京：三联书店，2016.
\item 郑永年.《中国模式》[M]. 北京：东方出版社，2010.
\item 韩庆祥. 中国特色社会主义理论体系研究[M]. 北京：人民出版社，2018.
\end{enumerate}

\end{document}
